\section{Conclusions}
To adapt to the sudden need for \gls{HALEU}, the \gls{DOE} is 
considering downblending some stockpiles of \gls{HEU} to produce 
fuel for reactors. However, the presence of other uranium isotopes 
in the fuel can introduce affects on the reactor performance that are 
not typically seen in fuel produced from enriching natural 
uranium. To investigate if these effects can significantly impact 
reactor performance, we simulated the reactor physics of 
an Xe-100-like reactor and an \gls{MMR}-like reactor using 
three different \gls{HALEU} compositions. The \gls{HALEU} 
compositions were based on pure \gls{HALEU} (only $^{235}$U and 
$^{238}$U), downblended \gls{HEU} from \gls{EBR} stockpiles, 
and downblended \gls{HEU} from the Y-12 National Security Complex. 
We compared the performance of the different fuel compositions
based on the \keff, \betaEff, energy- and spatially-dependent 
neutron flux, and the fuel, moderator, coolant, and total reactivity 
temperature feedback coefficients. 

In the Xe-100-like model, the impure fuel compositions resulted 
in higher \keff values, which suggests that additional neutron 
poisons may be needed during operation to control the reactivity 
in the core. The \betaEff values when using the impure fuel are 
lower than the \betaEff from the pure fuel, which suggests that 
reactivity control mechanisms may not be as effective when using 
the impure fuels. The impure fuels resulted in a depression in the 
thermal neutron flux across the active region of the core. There is 
is also an asymmetry in the axial fluxes that is partially because 
of the different fuel compositions and partially because of 
pebble placement in the core. Finally, the impure fuels 
lead to reactivity feedback coefficients that are more negative 
or within error of the values from the pure fuel. This result 
suggests that these fuel compositions does not significanty effect 
these metrics or would place the reactor in a safer operating state. 

In the \gls{MMR}-like model, the impure fuels result in smaller 
\keff values at each burnup step, but do not prevent the reactor 
from still being critical at \gls{EOL}. While this result may effect 
how neutron poisons are used during operation, it shows that 
the 20 year cycle time of the reactor is still possible with 
any of these \gls{HALEU} compositions. The \betaEff values from 
the different \gls{HALEU} compositions are all within error 
of each other at each burnup step, suggesting that the fuel 
composition does not significantly impact this metric for this 
reactor design. The impure fuels lead to larger fluxes in the 
fast energy range,, compared with the pure fuel, and the 
radially-dependent flux is dependent on proximity to fuel pin 
and coolant channels in the core. Most of the reactivity feedback 
coefficients from the different fuel compositions are within 
error of each other. The one exception is the coolant feedback 
coefficient, but the smaller absolute value of this coefficient 
compared with the others means that it has a smaller effect on 
the temperature-dependent reactivity changes in the core. These
results suggest that these different \gls{HALEU} compositions 
do not have a significant effect on the performance of this 
reactor design, and would not prevent it from operating in 
a safe state. 



\subsection{Future Work}
The work presented here is an initial start to understanding the impacts 
of using downblended \gls{HEU} in reactors. To more completely understand 
the effects of using these fuel compositions, one would need to 
update the models to more accurately capture each reactor design. This 
includes the reactor geometries, using non-homogeneous fuel materials, 
and nonhomogeneous material temperatures. Each of these model components 
limit the accuracy of the work completed here, and accurate depictions 
of them will provide more detail on the effects of using each fuel composition. 

Additionally, the Xe-100-like model assumes an equilibrium state. 
However, downblended \gls{HEU} may be used to create fuel for start up 
cores. Therefore, analysis of these compositions in the the Xe-100 
design in a start up core configuration would be more beneficial to 
understanding potential limitations of these fuel composisions 
to initially fuel reactors while developing infrastructure to 
enrich natural uranium to produce \gls{HALEU}.

Finally, neither of the companies designing these reactors have 
announced intentions to use downblended \gls{HEU}, but Oklo has announced 
intentions of using the downblended \gls{EBR} fuel for their Aurora reactor 
design \textbf{Find citation}. Therefore, a more practical investigation 
would 
be to apply this methodology to that reactor design, especially 
since it has some notable design differences to the reactors considered 
here, namely that it is a fast spectrum reactor. 
